\chapter*{When Things Go Wrong}
\markboth{When Things Go Wrong}{When Things Go Wrong}
\phantomsection
\addcontentsline{toc}{chapter}{When Things Go Wrong}

Every programmer spends time fixing programs. That is not the bad part of
programming. That is programming.

\section*{The Program Stops With An Error}
\phantomsection
\addcontentsline{toc}{section}{The Program Stops With An Error}

Read the line number first. Then list a few lines around it:

\begin{lstlisting}
LIST 80-120
\end{lstlisting}

Most beginner bugs are small: a missing quote, a wrong variable, a line number
that points to the wrong place, or a \texttt{NEXT} without the right
\texttt{FOR}.

\section*{The Program Runs But Looks Wrong}
\phantomsection
\addcontentsline{toc}{section}{The Program Runs But Looks Wrong}

If the program runs but the screen looks strange, check the setup lines. Look
for \texttt{MODE}, \texttt{COLOR}, \texttt{GCOLOR}, \texttt{GCLS}, and
\texttt{CLS}. A game that draws on the graphics screen needs the right mode and
the right colors before it starts drawing.

\section*{The Controls Do Nothing}
\phantomsection
\addcontentsline{toc}{section}{The Controls Do Nothing}

Look for the \texttt{GET} line. A typical key loop looks like this:

\begin{lstlisting}
90 GET K:IF K=0 THEN 90
\end{lstlisting}

That means ``keep waiting until a key appears.'' If the program waits forever,
check which key codes it accepts after that line.

\section*{The Best Debugging Tool}
\phantomsection
\addcontentsline{toc}{section}{The Best Debugging Tool}

Use \texttt{PRINT}. If you do not know what a variable contains, print it. If
you do not know whether a line is being reached, print a short message there.

\begin{lstlisting}
210 PRINT "X=";X;" Y=";Y
\end{lstlisting}

You can remove the extra \texttt{PRINT} lines after the program works.
